\FigureLayoutDeclare{img/2008/syllabus_paper_1_science/q12_fig1.png}{3.4cm}{17bp 17bp 17bp 4bp}

\chapter{2008年大纲I卷（理）}

\section{单选题}

\begin{problem}
函数 \( y = \sqrt{x(x - 1)} + \sqrt{x} \) 的定义域为（\quad）

\choices
    {\( \{x \mid x \geqslant 0\} \)}
    {\(\{x\mid x\geqslant 1\}\)}
    {\(\{x\mid x\geqslant 1\} \cup \{0\}\)}
    {\(\{x\mid 0\leqslant x\leqslant 1\}\)}
\end{problem}

\begin{answer}
C
\end{answer}

\begin{solution}
要使根式有意义，必须同时满足 \(x(x-1)\geqslant 0\) 和 \(x\geqslant 0\). 由 \(x(x-1)\geqslant 0\) 得 \(x\leqslant 0\) 或 \(x\geqslant 1\)，与 \(x\geqslant 0\) 取交集，得到 \(x=0\) 或 \(x\geqslant 1\). 所以定义域为 \(\{x\mid x\geqslant 1\}\cup\{0\}\)，故选 C.
\end{solution}

\begin{problem}
汽车经过启动，加速行驶，匀速行驶，减速行驶之后停车，若把这一过程中汽车的行驶路程 \(s\) 看作时间 \(t\) 的函数，其图象可能是（\quad）

\choices
    {\choicebitmap[width=0.15\paperwidth]{img/2008/syllabus_paper_1_science/q2_optA.png}}
    {\choicebitmap[width=0.15\paperwidth]{img/2008/syllabus_paper_1_science/q2_optB.png}}
    {\choicebitmap[width=0.15\paperwidth]{img/2008/syllabus_paper_1_science/q2_optC.png}}
    {\choicebitmap[width=0.15\paperwidth]{img/2008/syllabus_paper_1_science/q2_optD.png}}
\end{problem}

\begin{answer}
A
\end{answer}

\begin{solution}
设汽车的瞬时速度为 \(v(t)=s'(t)\). 启动后加速行驶时，\(v(t)\) 逐渐增大，故 \(s(t)\) 的斜率逐渐增大；匀速行驶时，\(v(t)\) 不变，故图象为斜率不变的上升直线；减速行驶时，\(v(t)\) 逐渐减小到 \(0\)，故图象上升但斜率逐渐减小，最后汽车停车时斜率为 \(0\). 四个图象中只有第一个图象同时具有上述三段变化特征，故选第一项.
\end{solution}

\begin{problem}
在 \(\triangle ABC\) 中，\(\overrightarrow{AB} = \bs{c}\)，\(\overrightarrow{AC} = \bs{b}\). 若点 \(D\) 满足 \(\overrightarrow{BD} = 2\overrightarrow{DC}\)，则 \(\overrightarrow{AD}=\)（\quad）

\choices
    {\(\frac{2}{3}\bs{b} + \frac{1}{3}\bs{c}\)}
    {\(\frac{5}{3}\bs{c} - \frac{2}{3}\bs{b}\)}
    {\(\frac{2}{3}\bs{b} - \frac{1}{3}\bs{c}\)}
    {\(\frac{1}{3}\bs{b} + \frac{2}{3}\bs{c}\)}
\end{problem}

\begin{answer}
A
\end{answer}

\begin{solution}
由 \(\overrightarrow{BD}=2\overrightarrow{DC}\) 可知 \(D\) 在线段 \(BC\) 上，且 \(BD:DC=2:1\)，所以 \(\overrightarrow{BD}=\frac{2}{3}\overrightarrow{BC}\). 又 \(\overrightarrow{BC}=\overrightarrow{AC}-\overrightarrow{AB}=\bs{b}-\bs{c}\)，因此
\[
\overrightarrow{AD}=\overrightarrow{AB}+\overrightarrow{BD}=\bs{c}+\frac{2}{3}(\bs{b}-\bs{c})=\frac{2}{3}\bs{b}+\frac{1}{3}\bs{c}.
\]
故选 A.
\end{solution}

\begin{problem}
设 \(a \in \R\)，且 \((a+\i)^2 \i\) 为正实数，则 \(a =\)（\quad）

\choices
    {\(2\)}
    {\(1\)}
    {\(0\)}
    {\(-1\)}
\end{problem}

\begin{answer}
D
\end{answer}

\begin{solution}
展开得
\[
(a+\i)^2\i=(a^2-1+2a\i)\i=-2a+(a^2-1)\i.
\]
它为正实数，故虚部为 \(0\)，实部为正数，即 \(a^2-1=0\) 且 \(-2a>0\). 由此 \(a=\pm1\)，再由 \(-2a>0\) 得 \(a=-1\)，故选 D.
\end{solution}

\begin{problem}
已知等差数列 \(\{a_{n}\}\) 满足 \(a_{2} + a_{4} = 4\)，\(a_{3} + a_{5} = 10\)，则它的前\(10\)项的和 \(S_{10} =\)（\quad）

\choices
    {\(138\)}
    {\(135\)}
    {\(95\)}
    {\(23\)}
\end{problem}

\begin{answer}
C
\end{answer}

\begin{solution}
设等差数列的首项为 \(a_1\)，公差为 \(d\). 由题意，
\[
a_2+a_4=2a_1+4d=4,\qquad a_3+a_5=2a_1+6d=10.
\]
两式相减得 \(2d=6\)，所以 \(d=3\)，再代入第一式得 \(a_1=-4\). 因而
\[
S_{10}=\frac{10}{2}\left(2a_1+9d\right)=5(-8+27)=95.
\]
故选 C.
\end{solution}

\begin{problem}
若函数 \( y = f(x - 1) \) 的图象与函数 \( y = \ln \sqrt{x} + 1 \) 的图象关于直线 \( y = x \) 对称，则 \( f(x) = \)（\quad）

\choices
    {\(\e^{2x - 1}\)}
    {\(\e^{2x}\)}
    {\(\e^{2x + 1}\)}
    {\(\e^{2x + 2}\)}
\end{problem}

\begin{answer}
B
\end{answer}

\begin{solution}
函数 \(y=\ln\sqrt{x}+1=\frac{1}{2}\ln x+1\)，其中 \(x>0\). 设 \(y=\frac{1}{2}\ln x+1\)，则 \(y-1=\frac{1}{2}\ln x\)，从而 \(x=\e^{2y-2}\). 交换 \(x,y\) 得其关于 \(y=x\) 的对称图象为 \(y=\e^{2x-2}\). 题设图象为 \(y=f(x-1)\)，故 \(f(x-1)=\e^{2x-2}\). 令 \(u=x-1\)，得到 \(f(u)=\e^{2u}\)，所以 \(f(x)=\e^{2x}\)，故选 B.
\end{solution}

\begin{problem}
设曲线 \( y = \frac{x + 1}{x - 1} \) 在点 \((3, 2)\) 处的切线与直线 \( ax + y + 1 = 0 \) 垂直，则 \( a = \)（\quad）

\choices
    {\(2\)}
    {\(\frac{1}{2}\)}
    {\(-\frac{1}{2}\)}
    {\(-2\)}
\end{problem}

\begin{answer}
D
\end{answer}

\begin{solution}
由 \(y=\frac{x+1}{x-1}\) 得
\[
y'=\frac{(x-1)-(x+1)}{(x-1)^2}=-\frac{2}{(x-1)^2}.
\]
所以曲线在 \((3,2)\) 处的切线斜率为 \(-\frac{1}{2}\). 直线 \(ax+y+1=0\) 的斜率为 \(-a\). 两直线垂直，故
\[
\left(-\frac{1}{2}\right)(-a)=-1,
\]
解得 \(a=-2\)，故选 D.
\end{solution}

\begin{problem}
为得到函数 \( y = \cos \left(2x + \frac{\pi}{3}\right) \) 的图象，只需将函数 \( y = \sin 2x \) 的图象（\quad）

\choices
    {向左平移 \(\frac{5\pi}{12}\) 个长度单位}
    {向右平移 \(\frac{5\pi}{12}\) 个长度单位}
    {向左平移 \(\frac{5\pi}{6}\) 个长度单位}
    {向右平移 \(\frac{5\pi}{6}\) 个长度单位}
\end{problem}

\begin{answer}
A
\end{answer}

\begin{solution}
利用诱导公式，
\[
\cos\left(2x+\frac{\pi}{3}\right)=\sin\left(\frac{\pi}{2}+2x+\frac{\pi}{3}\right)=\sin\left(2x+\frac{5\pi}{6}\right).
\]
将 \(y=\sin 2x\) 的图象向左平移 \(h\) 个长度单位后，函数变为 \(y=\sin 2(x+h)\). 令 \(2h=\frac{5\pi}{6}\)，得 \(h=\frac{5\pi}{12}\)，故选 A.
\end{solution}

\begin{problem}
设奇函数 \( f(x) \) 在 \((0, +\infty)\) 上为增函数，且 \( f(1) = 0 \)，则不等式 \( \frac{f(x) - f(-x)}{x} < 0 \) 的解集为（\quad）

\choices
    {\((-1,0)\cup (1, + \infty)\)}
    {\((- \infty, -1) \cup (0, 1)\)}
    {\((-\infty, -1) \cup (1, +\infty)\)}
    {\((-1,0)\cup (0,1)\)}
\end{problem}

\begin{answer}
D
\end{answer}

\begin{solution}
奇函数满足 \(f(-x)=-f(x)\)，所以原不等式等价于 \(\frac{2f(x)}{x}<0\)，且 \(x\ne 0\). 由 \(f\) 在 \((0,+\infty)\) 上递增且 \(f(1)=0\)，可得 \(f(x)<0\)（\(0<x<1\)），\(f(x)>0\)（\(x>1\)）. 当 \(-1<x<0\) 时，由奇性得 \(f(x)>0\)，故 \(\frac{f(x)}{x}<0\)；当 \(x<-1\) 时，\(f(x)<0\)，故 \(\frac{f(x)}{x}>0\). 综合得解集为 \((-1,0)\cup(0,1)\)，故选 D.
\end{solution}

\begin{problem}
若直线 \(\frac{x}{a} + \frac{y}{b} = 1\) 通过点 \(M(\cos \alpha, \sin \alpha)\)，则（\quad）

\choices
    {\(a^2 + b^2 \leqslant 1\)}
    {\(a^2 + b^2 \geqslant 1\)}
    {\(\frac{1}{a^2} + \frac{1}{b^2} \leqslant 1\)}
    {\(\frac{1}{a^2} + \frac{1}{b^2} \geqslant 1\)}
\end{problem}

\begin{answer}
D
\end{answer}

\begin{solution}
将 \(M(\cos\alpha,\sin\alpha)\) 代入直线方程，得
\[
\frac{\cos\alpha}{a}+\frac{\sin\alpha}{b}=1.
\]
由柯西不等式，
\[
1=\left(\frac{\cos\alpha}{a}+\frac{\sin\alpha}{b}\right)^2
\leqslant(\cos^2\alpha+\sin^2\alpha)\left(\frac{1}{a^2}+\frac{1}{b^2}\right)
=\frac{1}{a^2}+\frac{1}{b^2}.
\]
因此 \(\frac{1}{a^2}+\frac{1}{b^2}\geqslant 1\)，故选 D.
\end{solution}

\begin{problem}
已知三棱柱 \(ABC - A_{1}B_{1}C_{1}\) 的侧棱与底面边长都相等，\(A_{1}\) 在底面 \(ABC\) 内的射影为 \(\triangle ABC\) 的中心，则 \(AB_{1}\) 与底面 \(ABC\) 所成角的正弦值等于（\quad）

\choices
    {\(\frac{1}{3}\)}
    {\(\frac{\sqrt{2}}{3}\)}
    {\(\frac{\sqrt{3}}{3}\)}
    {\(\frac{2}{3}\)}
\end{problem}

\begin{answer}
B
\end{answer}

\begin{solution}
设等边三角形的边长为 \(a\)，建立空间直角坐标系，使
\[
A=(0,0,0),\quad B=(a,0,0),\quad C=\left(\frac{a}{2},\frac{\sqrt{3}a}{2},0\right).
\]
设 \(\triangle ABC\) 的中心为 \(O\)，则
\[
O=\left(\frac{a}{2},\frac{\sqrt{3}a}{6},0\right).
\]
由题意，\(A_1\) 在底面内的射影为 \(O\)，故可设 \(A_1=(\frac{a}{2},\frac{\sqrt{3}a}{6},h)\). 侧棱长等于 \(a\)，于是
\[
AA_1^2=\frac{a^2}{4}+\frac{a^2}{12}+h^2=a^2,\qquad h^2=\frac{2a^2}{3}.
\]
三棱柱的侧棱互相平行且相等，所以
\[
B_1=B+(A_1-A)=\left(\frac{3a}{2},\frac{\sqrt{3}a}{6},h\right).
\]
因此
\[
AB_1^2=\frac{9a^2}{4}+\frac{a^2}{12}+\frac{2a^2}{3}=3a^2.
\]
直线 \(AB_1\) 与底面所成角的正弦等于其竖直方向分量与长度之比，故
\[
\sin\theta=\frac{h}{AB_1}=\frac{\sqrt{2}a/\sqrt{3}}{\sqrt{3}a}=\frac{\sqrt{2}}{3}.
\]
故选 B.
\end{solution}

\begin{problem}
如图，一环形花坛分成 A，B，C，D 四块，现有 \(4\) 种不同的花供选种，要求在每块里种 \(1\) 种花，且相邻的 \(2\) 块种不同的花，则不同的种法总数为（\quad）

\bitmapfigure[max width=0.38\linewidth]{img/2008/syllabus_paper_1_science/q12_fig1.png}

\choices
    {\(96\)}
    {\(84\)}
    {\(60\)}
    {\(48\)}
\end{problem}

\begin{answer}
B
\end{answer}

\begin{solution}
先给区域 \(A\) 选花，有 \(4\) 种选法；给相邻区域 \(B\) 选花，有 \(3\) 种选法. 若 \(C\) 与 \(A\) 同色，则 \(C\) 只有 \(1\) 种选法，\(D\) 有 \(3\) 种选法；若 \(C\) 与 \(A\) 不同色，则 \(C\) 有 \(2\) 种选法，\(D\) 需同时不同于 \(A,C\)，有 \(2\) 种选法. 所以种法总数为
\[
4\times 3\times(1\times 3+2\times 2)=84.
\]
故选 B.
\end{solution}

\section{填空题}

\begin{problem}
若 \(x\)，\(y\) 满足约束条件 \(\begin{cases}
x + y \geqslant 0,\\
x - y + 3 \geqslant 0,\\
0 \leqslant x \leqslant 3,
\end{cases}\) 则 \(z = 2x - y\) 的最大值为\fillinblank{}.
\end{problem}

\begin{answer}
9
\end{answer}

\begin{solution}
由约束条件 \(x+y\geqslant 0\) 得 \(y\geqslant -x\). 对于固定的 \(x\)，\(z=2x-y\) 随 \(y\) 减小而增大，因此取 \(y=-x\) 时 \(z\) 最大. 此时 \(z=2x-(-x)=3x\)，又 \(0\leqslant x\leqslant 3\)，故最大值在 \(x=3\) 时取得，且 \(y=-3\) 满足其余约束条件. 所以 \(z_{\max}=9\).
\end{solution}

\begin{problem}
已知抛物线 \( y = ax^2 - 1 \) 的焦点是坐标原点，则以抛物线与两坐标轴的三个交点为顶点的三角形面积为\fillinblank{}.
\end{problem}

\begin{answer}
2
\end{answer}

\begin{solution}
将抛物线方程改写为 \(x^2=\frac{1}{a}(y+1)\). 与标准方程 \(x^2=4p(y+1)\) 比较，得 \(4p=\frac{1}{a}\). 抛物线顶点为 \((0,-1)\)，焦点为 \((0,-1+p)\)，题设焦点是原点，故 \(p=1\)，从而 \(a=\frac{1}{4}\). 因此抛物线与两坐标轴的交点为 \((0,-1)\)、\((-2,0)\)、\((2,0)\)，所成三角形的底为 \(4\)，高为 \(1\)，面积为 \(\frac{1}{2}\times4\times1=2\).
\end{solution}

\begin{problem}
在 \(\triangle ABC\) 中，\(AB = BC\)，\(\cos B = -\frac{7}{18}\). 若以 \(A\)，\(B\) 为焦点的椭圆经过点 \(C\)，则该椭圆的离心率 \(e = \)\fillinblank{}.
\end{problem}

\begin{answer}
\(\frac{3}{8}\)
\end{answer}

\begin{solution}
设 \(AB=BC=l\). 由余弦定理，
\[
AC^2=AB^2+BC^2-2AB\cdot BC\cos B
=2l^2\left(1+\frac{7}{18}\right)=\frac{25l^2}{9},
\]
所以 \(AC=\frac{5l}{3}\). 设椭圆的半长轴为 \(A_0\)，半焦距为 \(c_0\). 由椭圆定义，\(CA+CB=2A_0\)，故
\[
2A_0=\frac{5l}{3}+l=\frac{8l}{3},\qquad A_0=\frac{4l}{3}.
\]
椭圆焦点为 \(A,B\)，所以 \(2c_0=AB=l\)，即 \(c_0=\frac{l}{2}\). 因此离心率为
\[
e=\frac{c_0}{A_0}=\frac{l/2}{4l/3}=\frac{3}{8}.
\]
\end{solution}

\begin{problem}
等边三角形 \(ABC\) 与正方形 \(ABDE\) 有一公共边 \(AB\)，二面角 \(C - AB - D\) 的余弦值为 \(\frac{\sqrt{3}}{3}\)，\(M\)，\(N\) 分别是 \(AC\)，\(BC\) 的中点，则 \(EM\)，\(AN\) 所成角的余弦值等于\fillinblank{}.
\end{problem}

\begin{answer}
\(\frac{1}{6}\)
\end{answer}

\begin{solution}
不妨设 \(AB=2\)，建立空间直角坐标系，使正方形 \(ABDE\) 位于 \(xy\) 平面内：
\[
A=(0,0,0),\quad B=(2,0,0),\quad D=(2,2,0),\quad E=(0,2,0).
\]
设 \(C=(1,u,v)\). 由 \(AC=BC=AB=2\)，有 \(u^2+v^2=3\). 平面 \(ABC\) 内垂直于 \(AB\) 的方向可取 \((0,u,v)\)，正方形内垂直于 \(AB\) 的方向为 \((0,1,0)\)，故二面角条件给出
\[
\frac{u}{\sqrt{u^2+v^2}}=\frac{\sqrt{3}}{3},
\]
从而 \(u=1\)，\(v^2=2\). 取 \(v=\sqrt{2}\) 即可，另一种取向所得余弦相同. 因此
\[
M=\left(\frac{1}{2},\frac{1}{2},\frac{\sqrt{2}}{2}\right),\qquad N=\left(\frac{3}{2},\frac{1}{2},\frac{\sqrt{2}}{2}\right).
\]
于是
\[
\overrightarrow{EM}=\left(\frac{1}{2},-\frac{3}{2},\frac{\sqrt{2}}{2}\right),\qquad \overrightarrow{AN}=\left(\frac{3}{2},\frac{1}{2},\frac{\sqrt{2}}{2}\right).
\]
计算得
\[
\overrightarrow{EM}\cdot\overrightarrow{AN}=\frac{1}{2},\qquad |\overrightarrow{EM}|=|\overrightarrow{AN}|=\sqrt{3}.
\]
所以 \(EM\)，\(AN\) 所成角的余弦值为
\[
\frac{\overrightarrow{EM}\cdot\overrightarrow{AN}}{|\overrightarrow{EM}||\overrightarrow{AN}|}=\frac{1/2}{3}=\frac{1}{6}.
\]
\end{solution}

\section{解答题}

\begin{problem}
设 \(\triangle ABC\) 的内角 \(A\)，\(B\)，\(C\) 所对的边长分别为 \(a\)，\(b\)，\(c\)，且 \(a \cos B - b \cos A = \frac{3}{5}c\).

\begin{enumerate}
\item 求 \(\tan A \cot B\) 的值；
\item 求 \(\tan (A - B)\) 的最大值.
\end{enumerate}
\end{problem}

\begin{solution}
\begin{enumerate}
\item 由正弦定理，存在正数 \(k\) 使 \(a=k\sin A\)，\(b=k\sin B\)，\(c=k\sin C\). 代入已知等式并约去 \(k\)，得
\[
\sin A\cos B-\sin B\cos A=\frac{3}{5}\sin C.
\]
又 \(C=\pi-(A+B)\)，所以 \(\sin C=\sin(A+B)\). 因而
\[
\sin A\cos B-\cos A\sin B=\frac{3}{5}\left(\sin A\cos B+\cos A\sin B\right),
\]
整理得 \(\sin A\cos B=4\cos A\sin B\). 由于 \(\sin A>0\)，\(\sin B>0\)，上式说明 \(\cos A\) 与 \(\cos B\) 同号；又 \(A+B<\pi\)，所以 \(A\)，\(B\) 均为锐角，从而可除以 \(\cos A\cos B\)，得到 \(\tan A=4\tan B\)，进而 \(\tan A\cot B=4\).
\item 令 \(t=\tan B>0\)，则 \(\tan A=4t\). 于是
\[
\tan(A-B)=\frac{\tan A-\tan B}{1+\tan A\tan B}=\frac{3t}{1+4t^2}.
\]
由 \(1+4t^2\geqslant 4t\)，得
\[
\tan(A-B)\leqslant\frac{3t}{4t}=\frac{3}{4}.
\]
当 \(t=\frac{1}{2}\) 时等号成立，此时 \(A+B=\frac{\pi}{2}\)，确能构成三角形. 所以 \(\tan(A-B)\) 的最大值为 \(\frac{3}{4}\).
\end{enumerate}
\end{solution}

\begin{problem}
四棱锥 \(A - BCDE\) 中，底面 \(BCDE\) 为矩形，侧面 \(ABC \perp\) 底面 \(BCDE\)，\(BC = 2\)，\(CD = \sqrt{2}\)，\(AB = AC\).

\begin{enumerate}
\item 证明：\(AD \perp CE\)；
\item 设 \(CE\) 与平面 \(ABE\) 所成的角为 \(45^{\circ}\)，求二面角 \(C - AD - E\) 的大小.
\end{enumerate}

\bitmapfigure[max width=0.42\linewidth]{img/2008/syllabus_paper_1_science/q18_fig1.png}
\end{problem}

\begin{solution}
\begin{enumerate}
\item 建立空间直角坐标系，使
\[
B=(0,0,0),\quad C=(2,0,0),\quad D=(2,\sqrt{2},0),\quad E=(0,\sqrt{2},0).
\]
因平面 \(ABC\) 垂直于底面且含有 \(BC\)，可设 \(A=(x,0,h)\). 由 \(AB=AC\)，有
\[
x^2+h^2=(x-2)^2+h^2,
\]
所以 \(x=1\). 于是 \(A=(1,0,h)\). 计算
\[
\overrightarrow{AD}=(1,\sqrt{2},-h),\qquad \overrightarrow{CE}=(-2,\sqrt{2},0),
\]
从而 \(\overrightarrow{AD}\cdot\overrightarrow{CE}=-2+2=0\)，故 \(AD\perp CE\).
\item 由 \(\overrightarrow{AB}=(-1,0,-h)\)，\(\overrightarrow{BE}=(0,\sqrt{2},0)\)，可取平面 \(ABE\) 的法向量
\[
\bs{n}=\overrightarrow{AB}\times\overrightarrow{BE}=(h\sqrt{2},0,-\sqrt{2}).
\]
设 \(CE\) 与平面 \(ABE\) 所成的角为 \(\phi\)，则
\[
\sin\phi=\frac{|\overrightarrow{CE}\cdot\bs{n}|}{|\overrightarrow{CE}||\bs{n}|}
=\frac{2h\sqrt{2}}{\sqrt{6}\sqrt{2h^2+2}}.
\]
代入 \(\phi=45^{\circ}\)，平方化简得
\[
\frac{2h^2}{3(h^2+1)}=\frac{1}{2},
\]
所以 \(h^2=3\). 取 \(h=\sqrt{3}\)，则
\[
\overrightarrow{AD}=(1,\sqrt{2},-\sqrt{3}),\quad
\overrightarrow{AC}=(1,0,-\sqrt{3}),\quad
\overrightarrow{AE}=(-1,\sqrt{2},-\sqrt{3}).
\]
平面 \(ACD\) 与平面 \(AED\) 的法向量可分别取
\[
\bs{n}_1=\overrightarrow{AD}\times\overrightarrow{AC}=(-\sqrt{6},0,-\sqrt{2}),\qquad
\bs{n}_2=\overrightarrow{AD}\times\overrightarrow{AE}=(0,2\sqrt{3},2\sqrt{2}).
\]
两法向量的内积为 \(0\)，故两平面垂直，二面角 \(C-AD-E\) 的大小为 \(90^{\circ}\).
\end{enumerate}
\end{solution}

\begin{problem}
已知函数 \(f(x) = x^3 + ax^2 + x + 1\)，\(a \in \R\).

\begin{enumerate}
\item 讨论函数 \( f(x) \) 的单调区间；
\item 设函数 \( f(x) \) 在区间 \( \left(-\frac{2}{3}, -\frac{1}{3}\right) \) 内是减函数，求 \( a \) 的取值范围.
\end{enumerate}
\end{problem}

\begin{solution}
\begin{enumerate}
\item 函数的定义域为 \(\R\)，且
\[
f'(x)=3x^2+2ax+1.
\]
当 \(a^2\leqslant 3\) 时，二次式 \(3x^2+2ax+1\geqslant 0\) 对一切 \(x\) 成立，因此 \(f(x)\) 在 \(\R\) 上为增函数.

当 \(a^2>3\) 时，令
\[
x_1=\frac{-a-\sqrt{a^2-3}}{3},\qquad x_2=\frac{-a+\sqrt{a^2-3}}{3},
\]
则 \(x_1<x_2\)，且
\[
f'(x)=3(x-x_1)(x-x_2).
\]
所以 \(f(x)\) 在 \((-\infty,x_1)\) 和 \((x_2,+\infty)\) 上为增函数，在 \((x_1,x_2)\) 上为减函数.
\item 要使 \(f(x)\) 在 \(\left(-\frac{2}{3},-\frac{1}{3}\right)\) 内为减函数，需要 \(f'(x)\leqslant 0\) 对该区间内一切 \(x\) 成立，即
\[
3x^2+2ax+1\leqslant 0.
\]
由于 \(x<0\)，这等价于
\[
a\geqslant-\frac{3x^2+1}{2x}.
\]
令 \(g(x)=-\frac{3x^2+1}{2x}=-\frac{3}{2}x-\frac{1}{2x}\). 在所考察区间内，\(g(x)\) 的上确界为 \(2\)，因为当 \(x\to-\frac{1}{3}^{-}\) 时 \(g(x)\to 2\)，且 \(g(x)<2\) 对区间内的 \(x\) 成立. 所以必须有 \(a\geqslant2\). 反之，当 \(a\geqslant2\) 时，
\[
f'(x)\leqslant 3x^2+4x+1=(3x+1)(x+1)<0
\]
对 \(-\frac{2}{3}<x<-\frac{1}{3}\) 成立，故条件也充分. 因此 \(a\) 的取值范围为 \([2,+\infty)\).
\end{enumerate}
\end{solution}

\begin{problem}
已知 \(5\) 只动物中有 \(1\) 只患有某种疾病，需要通过化验血液来确定患病的动物. 血液化验结果呈阳性的即为患病动物，呈阴性即没患病. 下面是两种化验方案：

方案甲：逐个化验，直到能确定患病动物为止.

方案乙：先任取 \(3\) 只，将它们的血液混在一起化验. 若结果呈阳性则表明患病动物为这 \(3\) 只中的 \(1\) 只，然后再逐个化验，直到能确定患病动物为止；若结果呈阴性则在另外 \(2\) 只中任取 \(1\) 只化验.

\begin{enumerate}
\item 求依方案甲所需化验次数不少于依方案乙所需化验次数的概率；
\item \(\xi\) 表示依方案乙所需化验次数，求 \(\xi\) 的期望.
\end{enumerate}
\end{problem}

\begin{solution}
\begin{enumerate}
\item 设方案甲所需化验次数为 \(X\)，方案乙所需化验次数为 \(Y\). 方案甲按固定顺序逐只化验，患病动物在第 \(1,2,3\) 个位置的概率均为 \(\frac{1}{5}\)，所以
\[
P(X\geqslant2)=1-\frac{1}{5}=\frac{4}{5},\qquad P(X\geqslant3)=1-\frac{1}{5}-\frac{1}{5}=\frac{3}{5}.
\]
对方案乙，先取的三只中有患病动物的概率为 \(\frac{3}{5}\). 在这种情况下，若逐个化验时第一次就验出患病动物，方案乙共化验 \(2\) 次，其概率为
\[
\frac{3}{5}\times\frac{1}{3}=\frac{1}{5};
\]
若第一次没有验出，则还需再化验一次，共化验 \(3\) 次，其概率为 \(\frac{3}{5}\times\frac{2}{3}=\frac{2}{5}\). 若混合化验呈阴性，概率为 \(\frac{2}{5}\)，再从另外两只中任取一只化验后即可确定患病动物，所以共化验 \(2\) 次. 因而
\[
P(Y=2)=\frac{1}{5}+\frac{2}{5}=\frac{3}{5},\qquad P(Y=3)=\frac{2}{5}.
\]
按两种方案的随机选择独立处理，得
\[
P(X\geqslant Y)=P(Y=2)P(X\geqslant2)+P(Y=3)P(X\geqslant3)
=\frac{3}{5}\times\frac{4}{5}+\frac{2}{5}\times\frac{3}{5}
=\frac{18}{25}.
\]
\item 由上面的分布，
\[
E(\xi)=2P(Y=2)+3P(Y=3)
=2\times\frac{3}{5}+3\times\frac{2}{5}
=\frac{12}{5}.
\]
\end{enumerate}
\end{solution}

\begin{problem}
双曲线的中心为原点 \(O\)，焦点在 \(x\) 轴上，两条渐近线分别为 \(l_{1}\)，\(l_{2}\)，经过右焦点 \(F\) 垂直于 \(l_{1}\) 的直线分别交 \(l_{1}\)，\(l_{2}\) 于 \(A\)，\(B\) 两点. 已知 \(\left|\overrightarrow{OA}\right|\)，\(\left|\overrightarrow{AB}\right|\)，\(\left|\overrightarrow{OB}\right|\) 成等差数列，且 \(\overrightarrow{BF}\) 与 \(\overrightarrow{FA}\) 同向.

\begin{enumerate}
\item 求双曲线的离心率；
\item 设 \(AB\) 被双曲线所截得的线段的长为 \(4\)，求双曲线的方程.
\end{enumerate}
\end{problem}

\begin{solution}
\begin{enumerate}
\item 设双曲线方程为 \(\frac{x^2}{a^2}-\frac{y^2}{b^2}=1\)，其中 \(a>0\)，\(b>0\)，焦距参数满足 \(c^2=a^2+b^2\). 令 \(k=\frac{b}{a}\)，则两条渐近线为 \(y=\pm kx\). 由 \(\overrightarrow{BF}\) 与 \(\overrightarrow{FA}\) 同向可知 \(F\) 在线段 \(BA\) 之间，从而 \(0<k<1\).

取 \(l_1:y=kx\)，过 \(F=(c,0)\) 作其垂线，方程为 \(y=-\frac{x-c}{k}\). 它与 \(l_1,l_2\) 的交点分别为
\[
A=\left(\frac{c}{1+k^2},\frac{kc}{1+k^2}\right),\qquad
B=\left(\frac{c}{1-k^2},-\frac{kc}{1-k^2}\right).
\]
由 \(c=a\sqrt{1+k^2}\)，可得
\[
|OA|=a,\qquad |OB|=\frac{a(1+k^2)}{1-k^2},\qquad |AB|=\frac{2ak}{1-k^2}.
\]
三者成等差数列，所以
\[
a+\frac{a(1+k^2)}{1-k^2}=2\cdot\frac{2ak}{1-k^2},
\]
化简得 \(k=\frac{1}{2}\). 因而离心率为
\[
e=\frac{c}{a}=\sqrt{1+k^2}=\frac{\sqrt{5}}{2}.
\]
\item 由 \(k=\frac{1}{2}\)，得 \(b=\frac{a}{2}\)，\(c=\frac{\sqrt{5}}{2}a\)，且直线 \(AB\) 的斜率为 \(-\frac{1}{k}=-2\). 因为它经过 \(F=(c,0)\)，故
\[
y=-2(x-c).
\]
将其代入双曲线方程 \(x^2-4y^2=a^2\)，得
\[
15x^2-32cx+21a^2=0.
\]
设双曲线与直线 \(AB\) 的两个交点横坐标为 \(x_1,x_2\)，则
\[
|x_1-x_2|
=\frac{\sqrt{(32c)^2-4\cdot15\cdot21a^2}}{15}
=\frac{2\sqrt{5}a}{15}.
\]
直线 \(AB\) 的斜率为 \(-2\)，所以弦长为
\[
\sqrt{1+(-2)^2}|x_1-x_2|=\sqrt{5}\cdot\frac{2\sqrt{5}a}{15}=\frac{2a}{3}.
\]
由弦长为 \(4\) 得 \(a=6\)，从而 \(b=3\). 所求双曲线方程为
\[
\frac{x^2}{36}-\frac{y^2}{9}=1.
\]
\end{enumerate}
\end{solution}

\begin{problem}
设函数 \( f(x) = x - x \ln x \). 数列 \( \{a_n\} \) 满足 \(0 < a_1 < 1\)，\(a_{n+1} = f(a_n)\).

\begin{enumerate}
\item 证明：函数 \( f(x) \) 在区间 \((0, 1)\) 是增函数；
\item 证明：\( a_{n}<a_{n+1}<1 \)；
\item 设 \( b \in (a_1, 1) \)，整数 \( k \geqslant \frac{a_1 - b}{a_1 \ln b} \). 证明：\( a_{k+1} > b \).
\end{enumerate}
\end{problem}

\begin{solution}
\begin{enumerate}
\item 在 \((0,1)\) 上，
\[
f'(x)=1-(\ln x+1)=-\ln x>0.
\]
所以 \(f(x)\) 在 \((0,1)\) 上为增函数.
\item 对 \(0<x<1\)，有
\[
f(x)-x=-x\ln x>0.
\]
又由于 \(f\) 在 \((0,1)\) 上递增，并且可连续延拓到 \(x=1\) 且 \(f(1)=1\)，故 \(f(x)<1\). 由 \(0<a_1<1\)，归纳可得：若 \(0<a_n<1\)，则
\[
a_{n+1}=f(a_n)>a_n,\qquad a_{n+1}<1.
\]
因此对一切正整数 \(n\)，都有 \(a_n<a_{n+1}<1\).
\item 由 \(a_1<b<1\) 可知
\[
\frac{a_1-b}{a_1\ln b}>0,
\]
又因 \(k\) 作为数列下标取正整数，故 \(k\geqslant1\). 反设 \(a_{k+1}\leqslant b\). 由上一问，\(a_n\) 递增，故 \(a_1\leqslant a_n\leqslant b\)（\(1\leqslant n\leqslant k+1\)），从而
\[
a_{n+1}-a_n=-a_n\ln a_n\geqslant-a_1\ln b.
\]
当 \(n=1\) 时因 \(a_1<b\)，不等式严格成立，所以
\[
a_{k+1}-a_1=\sum_{n=1}^{k}(a_{n+1}-a_n)>-ka_1\ln b.
\]
由 \(k\geqslant\frac{a_1-b}{a_1\ln b}\) 且 \(a_1\ln b<0\)，不等式乘以 \(-a_1\ln b>0\) 后得
\[
-ka_1\ln b\geqslant b-a_1.
\]
于是 \(a_{k+1}>a_1+b-a_1=b\)，与反设矛盾. 因此 \(a_{k+1}>b\).
\end{enumerate}
\end{solution}
